% Options for packages loaded elsewhere
\PassOptionsToPackage{unicode}{hyperref}
\PassOptionsToPackage{hyphens}{url}
\PassOptionsToPackage{dvipsnames,svgnames,x11names}{xcolor}
%
\documentclass[
  11pt,
]{article}
\usepackage{amsmath,amssymb}
\usepackage{iftex}
\ifPDFTeX
  \usepackage[T1]{fontenc}
  \usepackage[utf8]{inputenc}
  \usepackage{textcomp} % provide euro and other symbols
\else % if luatex or xetex
  \usepackage{unicode-math} % this also loads fontspec
  \defaultfontfeatures{Scale=MatchLowercase}
  \defaultfontfeatures[\rmfamily]{Ligatures=TeX,Scale=1}
\fi
\usepackage{lmodern}
\ifPDFTeX\else
  % xetex/luatex font selection
\fi
% Use upquote if available, for straight quotes in verbatim environments
\IfFileExists{upquote.sty}{\usepackage{upquote}}{}
\IfFileExists{microtype.sty}{% use microtype if available
  \usepackage[]{microtype}
  \UseMicrotypeSet[protrusion]{basicmath} % disable protrusion for tt fonts
}{}
\makeatletter
\@ifundefined{KOMAClassName}{% if non-KOMA class
  \IfFileExists{parskip.sty}{%
    \usepackage{parskip}
  }{% else
    \setlength{\parindent}{0pt}
    \setlength{\parskip}{6pt plus 2pt minus 1pt}}
}{% if KOMA class
  \KOMAoptions{parskip=half}}
\makeatother
\usepackage{xcolor}
\usepackage{microtype}
\usepackage[a4paper,margin=1in]{geometry}
\usepackage{color}
\usepackage{fancyvrb}
\newcommand{\VerbBar}{|}
\newcommand{\VERB}{\Verb[commandchars=\\\{\}]}
\DefineVerbatimEnvironment{Highlighting}{Verbatim}{commandchars=\\\{\},fontsize=\small}
% Add ',fontsize=\small' for more characters per line
\newenvironment{Shaded}{}{}
\newcommand{\AlertTok}[1]{\textcolor[rgb]{1.00,0.00,0.00}{\textbf{#1}}}
\newcommand{\AnnotationTok}[1]{\textcolor[rgb]{0.38,0.63,0.69}{\textbf{\textit{#1}}}}
\newcommand{\AttributeTok}[1]{\textcolor[rgb]{0.49,0.56,0.16}{#1}}
\newcommand{\BaseNTok}[1]{\textcolor[rgb]{0.25,0.63,0.44}{#1}}
\newcommand{\BuiltInTok}[1]{\textcolor[rgb]{0.00,0.50,0.00}{#1}}
\newcommand{\CharTok}[1]{\textcolor[rgb]{0.25,0.44,0.63}{#1}}
\newcommand{\CommentTok}[1]{\textcolor[rgb]{0.38,0.63,0.69}{\textit{#1}}}
\newcommand{\CommentVarTok}[1]{\textcolor[rgb]{0.38,0.63,0.69}{\textbf{\textit{#1}}}}
\newcommand{\ConstantTok}[1]{\textcolor[rgb]{0.53,0.00,0.00}{#1}}
\newcommand{\ControlFlowTok}[1]{\textcolor[rgb]{0.00,0.44,0.13}{\textbf{#1}}}
\newcommand{\DataTypeTok}[1]{\textcolor[rgb]{0.56,0.13,0.00}{#1}}
\newcommand{\DecValTok}[1]{\textcolor[rgb]{0.25,0.63,0.44}{#1}}
\newcommand{\DocumentationTok}[1]{\textcolor[rgb]{0.73,0.13,0.13}{\textit{#1}}}
\newcommand{\ErrorTok}[1]{\textcolor[rgb]{1.00,0.00,0.00}{\textbf{#1}}}
\newcommand{\ExtensionTok}[1]{#1}
\newcommand{\FloatTok}[1]{\textcolor[rgb]{0.25,0.63,0.44}{#1}}
\newcommand{\FunctionTok}[1]{\textcolor[rgb]{0.02,0.16,0.49}{#1}}
\newcommand{\ImportTok}[1]{\textcolor[rgb]{0.00,0.50,0.00}{\textbf{#1}}}
\newcommand{\InformationTok}[1]{\textcolor[rgb]{0.38,0.63,0.69}{\textbf{\textit{#1}}}}
\newcommand{\KeywordTok}[1]{\textcolor[rgb]{0.00,0.44,0.13}{\textbf{#1}}}
\newcommand{\NormalTok}[1]{#1}
\newcommand{\OperatorTok}[1]{\textcolor[rgb]{0.40,0.40,0.40}{#1}}
\newcommand{\OtherTok}[1]{\textcolor[rgb]{0.00,0.44,0.13}{#1}}
\newcommand{\PreprocessorTok}[1]{\textcolor[rgb]{0.74,0.48,0.00}{#1}}
\newcommand{\RegionMarkerTok}[1]{#1}
\newcommand{\SpecialCharTok}[1]{\textcolor[rgb]{0.25,0.44,0.63}{#1}}
\newcommand{\SpecialStringTok}[1]{\textcolor[rgb]{0.73,0.40,0.53}{#1}}
\newcommand{\StringTok}[1]{\textcolor[rgb]{0.25,0.44,0.63}{#1}}
\newcommand{\VariableTok}[1]{\textcolor[rgb]{0.10,0.09,0.49}{#1}}
\newcommand{\VerbatimStringTok}[1]{\textcolor[rgb]{0.25,0.44,0.63}{#1}}
\newcommand{\WarningTok}[1]{\textcolor[rgb]{0.38,0.63,0.69}{\textbf{\textit{#1}}}}
\usepackage{longtable,booktabs,array}
\usepackage{calc} % for calculating minipage widths
% Correct order of tables after \paragraph or \subparagraph
\usepackage{etoolbox}
\makeatletter
\patchcmd\longtable{\par}{\if@noskipsec\mbox{}\fi\par}{}{}
\makeatother
% Allow footnotes in longtable head/foot
\IfFileExists{footnotehyper.sty}{\usepackage{footnotehyper}}{\usepackage{footnote}}
\makesavenoteenv{longtable}
\usepackage{graphicx}
\usepackage{float}
\usepackage{caption}
\makeatletter
\def\maxwidth{\ifdim\Gin@nat@width>\linewidth\linewidth\else\Gin@nat@width\fi}
\def\maxheight{\ifdim\Gin@nat@height>\textheight\textheight\else\Gin@nat@height\fi}
\makeatother
% Scale images if necessary, so that they will not overflow the page
% margins by default, and it is still possible to overwrite the defaults
% using explicit options in \includegraphics[width, height, ...]{}
\setkeys{Gin}{width=\maxwidth,height=\maxheight,keepaspectratio}
% Set default figure placement to htbp
\makeatletter
\def\fps@figure{htbp}
\makeatother
\setlength{\emergencystretch}{8em} % prevent overfull lines
\providecommand{\tightlist}{%
  \setlength{\itemsep}{0pt}\setlength{\parskip}{0pt}}
\setcounter{secnumdepth}{-\maxdimen} % remove section numbering
\ifLuaTeX
\usepackage[bidi=basic]{babel}
\else
\usepackage[bidi=default]{babel}
\fi
\babelprovide[main,import]{american}
% get rid of language-specific shorthands (see #6817):
\let\LanguageShortHands\languageshorthands
\def\languageshorthands#1{}
\ifLuaTeX
  \usepackage{selnolig}  % disable illegal ligatures
\fi
\usepackage{bookmark}
\IfFileExists{xurl.sty}{\usepackage{xurl}}{} % add URL line breaks if available
\urlstyle{same}
\hypersetup{
  pdftitle={Linear-Error Clique Partitions of Split Graphs via Structured Triangle Packing},
  pdflang={en-US},
  colorlinks=true,
  linkcolor={blue},
  filecolor={Maroon},
  citecolor={Blue},
  urlcolor={Blue},
  pdfcreator={LaTeX via pandoc}}

\title{Linear-Error Clique Partitions of Split Graphs via Structured Triangle Packing}
\author{Juan Pablo Traverso Gianini\\
\small Independent researcher, Santiago, Chile\\
\small \href{mailto:jtraverso@gmail.com}{jtraverso@gmail.com}\\
\small ORCID: 0009-0003-6068-4096}
\date{July 21, 2026}

\begin{document}
\sloppy
\maketitle
\begin{small}\begin{raggedright}

\textbf{Paper III in the series}\\
\textbf{Author manuscript:} v0.9.5 style-matched review draft\\
\textbf{Date:} July 21, 2026\\
\textbf{Status:} working manuscript; formal verification in progress. The Lean 4 / Mathlib v4.28.0 project currently builds as a complete dependency graph. Its elementary core is reduced to eight localized proof obligations, with no \texttt{admit}, \texttt{native\_decide}, \texttt{unsafe}, or unrecorded project axiom. The only intended mathematical axioms are the two external asymptotic inputs stated in Section 2.4. The release wording ``formally verified'' and ``sorry-free'' is reserved until the final build, axiom report, and commit are frozen. Not externally peer-reviewed; bibliographic and novelty review incomplete. \textbf{MSC 2020:} Primary 05C70; Secondary 05C35, 05C72, 05C15.

\begin{center}\rule{0.5\linewidth}{0.5pt}\end{center}

\end{raggedright}\end{small}

\section{Abstract}\label{abstract}

Let \(G=(K\cup I,E)\) be a split graph on \(n\) vertices, with \(K\) a
clique and \(I\) an independent set. We prove that there is an absolute
constant \(C\) such that

\[
|E(G)|-2\nu_3(G)
\le
\frac{n^2}{6}+Cn,
\]

where \(\nu_3(G)\) is the maximum number of pairwise edge-disjoint
triangles in \(G\). Consequently,

\[
\operatorname{cp}(G)
\le
\frac{n^2}{6}+Cn
\]

for every split graph \(G\), where \(\operatorname{cp}(G)\) is the
minimum number of cliques whose edge sets partition \(E(G)\). This
strengthens the asymptotic \((1/6+o(1))n^2\) bound to the correct
linear-error scale. The leading constant is sharp.

The proof separates three regimes according to \(\alpha=|I|/|K|\). In
the bulk regime, an exact four-orbit linear program for a
common-neighborhood profile yields a uniform quadratic fractional
margin; the Haxell--Rödl theorem then absorbs the subquadratic
integrality loss. When \(\alpha\to0\), large edge-disjoint matchings
anchored at the independent vertices leave an almost complete,
triangle-divisible clique residual, which is decomposed exactly by dense
graph decomposition results. When \(\alpha\to2\), averaged factorization
closes the short corridor, while a double-factor polarization inequality
and a shifted-center gain-completion argument close the remaining
mesoscopic corridor.

The proof does not establish the stronger universal estimate
\(\nu_3^*(G)-\nu_3(G)=O(n)\). Away from the extremal corridor, a
potentially superlinear integrality gap is absorbed by quadratic
fractional slack.

All corridor mechanisms are self-contained; in particular, the list edge
coloring ingredient is proved from first principles in Appendix D. Only
two external asymptotic theorems remain, confined to the bulk and sparse
regimes, and the near-extremal corridor theorem (Proposition 10.1) is
fully elementary and effective.

\textbf{Keywords:} clique partition; split graph; triangle packing;
fractional packing; factorization; polarization; list edge coloring;
graph decomposition.

\begin{center}\rule{0.5\linewidth}{0.5pt}\end{center}

\section{1. Introduction}\label{introduction}

\subsection{1.1 Clique partitions of split
graphs}\label{clique-partitions-of-split-graphs}

A \textbf{clique partition} of a graph \(G\) is a family of complete
subgraphs whose edge sets partition \(E(G)\). Cliques of order two are
allowed. The minimum size of such a family is denoted by
\(\operatorname{cp}(G)\).

The chordal clique-partition problem of Erdős, Ordman, and Zalcstein
asks for the correct extremal order of \(\operatorname{cp}(G)\) on
chordal graphs. Their complete-split (indeed threshold) construction
shows that the coefficient \(1/6\) is unavoidable. Split graphs form a
natural and important subclass of chordal graphs: a graph is split if
its vertex set can be partitioned into a clique and an independent set.

This is the third paper in the series on the Erdős--Ordman--Zalcstein
clique-partition problem (Erdős Problem \#81). Paper I proves a finite
fractional split-graph bound, \(|E(G)|-2\nu_3^*(G)\le n^2/6+n\). Paper
II identifies the complete-split terminal family and determines the
exact finite maximum of the fractional cover functional over chordal
graphs. Papers I and II are fractional and are formally verified in Lean
4. The present paper is the first integral installment: it strengthens
the split-graph estimate to the linear-error scale \(n^2/6+O(n)\).

Paper II of the series determines the fractional extremal problem
exactly. For every \(n\),

\[
\max_{\substack{|V(G)|=n\\G\text{ chordal}}}
\bigl(|E(G)|-2\nu_3^*(G)\bigr)
=
\max_{\substack{|V(G)|=n\\G\text{ split}}}
\bigl(|E(G)|-2\nu_3^*(G)\bigr)
=
\left\lfloor\frac{(2n+1)^2}{24}\right\rfloor.
\]

It follows, by standard fractional-to-integral rounding, that every
chordal graph satisfies

\[
\operatorname{cp}(G)
\le
\left(\frac16+o(1)\right)n^2.
\]

The present paper addresses a different and genuinely second-order
question: it strengthens the integral bound to \(n^2/6+O(n)\) for
arbitrary split graphs.

\subsection{1.2 Triangle packings and the linear-error
target}\label{triangle-packings-and-the-linear-error-target}

Let \(\nu_3(G)\) be the maximum number of pairwise edge-disjoint
triangles in \(G\). Every such packing gives a clique partition: keep
every packed triangle and use one \(K_2\) for each uncovered edge. Hence

\[
\operatorname{cp}(G)
\le
|E(G)|-2\nu_3(G).
\tag{1.1}
\]

The central theorem of this paper is an integral estimate at the correct
second-order scale.

\subsubsection{Theorem 1.1}\label{theorem-1.1-linear-error-split-graph-theorem}

There exists an absolute constant \(C\) such that every split graph
\(G\) on \(n\) vertices satisfies

\[
\boxed{
|E(G)|-2\nu_3(G)
\le
\frac{n^2}{6}+Cn.
}
\tag{1.2}
\]

Combining (1.1) and Theorem 1.1 gives the clique-partition statement.

\subsubsection{Corollary 1.2}\label{corollary-1.2-linear-error-clique-partition-bound}

There exists an absolute constant \(C\) such that every split graph
\(G\) on \(n\) vertices satisfies

\[
\boxed{
\operatorname{cp}(G)
\le
\frac{n^2}{6}+Cn.
}
\tag{1.3}
\]

The Erdős--Ordman--Zalcstein complete-split (indeed threshold) family
\(K_p\vee\overline K_{2p}\), with \(n=3p\), shows that the leading
constant \(1/6\) is sharp and that a linear term naturally appears.

\begin{figure}[H]
\centering
\includegraphics[width=0.72\textwidth]{paperIII_figures_bilingual_v0.9.7/fig1_complete_split_extremizer_en.png}
\caption{The complete-split graph \(K_p\vee\overline K_{2p}\). Every edge between the two parts is present, and the independent side has no internal edges. Since every triangle contains a clique edge, an edge-disjoint triangle packing uses at most \(\binom p2\) triangles; see Section 10.2. The figure is illustrative and is not used as a premise of the proof.}
\label{fig:complete-split}
\end{figure}

\subsection{1.3 Localization of hypothetical
obstructions}\label{localization-of-hypothetical-obstructions}

\subsubsection{Proposition 1.3}\label{proposition-1.3-localization-of-hypothetical-obstructions}

Suppose, for contradiction, that no absolute linear-error bound holds,
and for each positive integer \(k\) choose a minimum-order split graph
\(G_k\) satisfying

\[
|E(G_k)|-2\nu_3(G_k)>\frac{|V(G_k)|^2}{6}+k|V(G_k)|.
\]

After passing to a subsequence, writing \(|K_k|=p_k\), \(|I_k|=q_k\),
every such hypothetical sequence must satisfy

\[
\frac{q_k}{p_k}\longrightarrow 2.
\]

Moreover, if \(q_k=2p_k-s_k\), then the bulk and short-corridor
estimates force

\[
\sqrt{p_k}\ll s_k=o(p_k).
\]

Thus any obstruction to a linear error term would have to concentrate in
the mesoscopic near-extremal corridor. The high/low-dispersion dichotomy
in Section 9 excludes that corridor.

\subsection{1.4 What the theorem does not
prove}\label{what-the-theorem-does-not-prove}

Let \(\nu_3^*(G)\) denote the fractional triangle-packing number. The
present proof does \textbf{not} establish

\[
\nu_3^*(G)-\nu_3(G)=O(n)
\tag{1.4}
\]

uniformly over split graphs.

Indeed, write

\[
\Delta(G)=\nu_3^*(G)-\nu_3(G)
\]

and

\[
S(G)=\frac{n^2}{6}-\bigl(|E(G)|-2\nu_3^*(G)\bigr).
\]

Then

\[
|E(G)|-2\nu_3(G)
=
\frac{n^2}{6}-S(G)+2\Delta(G).
\tag{1.5}
\]

Theorem 1.1 only requires

\[
2\Delta(G)\le S(G)+O(n).
\]

In the bulk regime, \(S(G)\) is quadratic and absorbs the general
\(o(n^2)\) integrality loss. Thus the universal linear integrality-gap
problem remains open.

\subsection{1.5 Proof architecture}\label{proof-architecture}

Write

\[
|K|=p,
\qquad
|I|=q,
\qquad
\alpha=\frac qp.
\]

The ratio \(\alpha=q/p\) determines which source of slack is available.
We first dispose of \(q\ge2p-1\) by direct averaged factorization, and
may then assume \(0\le\alpha<2\). A hypothetical sequence of
counterexamples has a subsequence in one of three regimes.

\begin{enumerate}
\def\labelenumi{\arabic{enumi}.}
\item
  \textbf{Bulk:} \(\alpha\) stays away from both \(0\) and \(2\). An
  exact common-profile LP and a fractional cloning argument produce
  quadratic fractional slack. Haxell--Rödl converts the fractional
  packing into an integral one with only \(o(n^2)\) loss.
\item
  \textbf{Sparse independent side:} \(\alpha\to0\). Each independent
  vertex anchors a large matching inside its neighborhood. The residual
  clique is almost complete. After deleting \(O(p)\) edges to correct
  divisibility, it admits an exact triangle decomposition.
\item
  \textbf{Near-extremal corridor:} \(\alpha\to2\). Write \(q=2p-s\).
  Averaged factorization closes \(s=O(\sqrt p)\). For
  \(\sqrt p\ll s=o(p)\), high dispersion is paid by a double-factor
  polarization term, while low dispersion implies closeness to a common
  center and is handled by shifted-center gain completion.
\end{enumerate}

The division is not cosmetic. In the bulk there is enough quadratic
fractional slack to absorb a general \(o(n^2)\) rounding loss. At the
two ends that slack vanishes, so the proof replaces generic rounding by
constructions adapted to the geometry of the split presentation.

\begin{figure}[H]
\centering
\includegraphics[width=0.94\textwidth]{paperIII_figures_bilingual_v0.9.7/fig2_alpha_regimes_en.png}
\caption{The three parameter regimes used in the proof. The principal corridor thresholds are shown; the complete hypotheses are stated in Sections 4, 8, and 9. The endpoint regimes are asymptotic rather than hard subintervals of \([0,2]\).}
\label{fig:alpha-regimes}
\end{figure}

\subsection{1.6 Supplementary computational
audits}\label{supplementary-computational-audits}

The accompanying closure package contains exact-arithmetic and
finite-instance regression tests for selected identities and
quantitative inequalities used in the proof.

Every audited statement is proved analytically in the manuscript. No
computation, script, finite enumeration, or solver output is a logical
premise of Theorem 1.1. The supplementary material is provided solely
for independent checking, regression testing, and reproducibility.

\subsection{1.7 Organization}\label{organization}

Section 2 fixes notation and external inputs. Section 3 solves the
common-profile LP. Section 4 proves the fractional cloning bound and the
global fractional margin. Section 5 develops averaged and double-factor
rounding. Section 6 proves polarization. Section 7 proves shifted-center
gain completion. Section 8 handles the sparse-independent regime.
Section 9 assembles the three regimes. Section 10 records corollaries,
Section 11 discusses the proof and its limitations, and Section 12
describes future uses and open directions. The appendices contain
algebraic details, the divisibility correction, reproducibility
material, and a self-contained proof of the list edge coloring theorem
(Appendix D).

\begin{center}\rule{0.5\linewidth}{0.5pt}\end{center}

\section{2. Preliminaries}\label{preliminaries}

\subsection{2.1 Split notation}\label{split-notation}

Throughout,

\[
G=(K\sqcup I,E)
\]

is a split graph, where \(K\) is a clique of size \(p\) and
\(I=\{v_1,\ldots,v_q\}\) is independent.

For \(v_i\in I\), write

\[
N_i=N(v_i)\cap K,
\qquad
d_i=|N_i|,
\qquad
S_i=K\setminus N_i,
\qquad
m_i=|S_i|.
\]

Set

\[
M=\sum_i m_i,
\qquad
S_2=\sum_i m_i^2.
\]

When \(q\) is close to \(2p\), write

\[
q=2p-s.
\tag{2.1}
\]

\subsection{2.2 Triangle packings}\label{triangle-packings}

A fractional triangle packing is a nonnegative weight assignment to the
triangles of \(G\) such that the total weight through every edge is at
most one. Its maximum value is \(\nu_3^*(G)\).

The dual is a fractional triangle cover: a nonnegative weight assignment
to edges such that every triangle receives total edge weight at least
one. LP duality gives equality between the two optimum values.

\subsection{2.3 Complete graph
factorizations}\label{complete-graph-factorizations}

Adopt the boundary conventions

\[
\chi'(K_0)=\chi'(K_1)=0.
\]

For \(t\ge2\), the edge-chromatic number of a complete graph is

\[
\chi'(K_t)
=
\begin{cases}
t-1,&t\text{ even},\\
t,&t\text{ odd}.
\end{cases}
\tag{2.2}
\]

Thus \(E(K_t)\) decomposes into \(\chi'(K_t)\) matchings. The
conventions for \(t\le1\) make the same language valid for empty edge
sets.

\subsection{2.4 External theorems}\label{external-theorems}

We use two external asymptotic inputs. (Earlier versions cited a third,
the list edge coloring theorem; the case required by this paper is
proved self-contained in Appendix D, so Sections 5--7 and Proposition
10.1 depend on no external theorem.)

\subsubsection{Theorem 2.1 ---
Haxell--Rödl/Yuster}\label{theorem-2.1-haxellruxf6dlyuster}

For every fixed graph \(H\),

\[
\nu_H^*(G)-\nu_H(G)=o(|V(G)|^2)
\]

uniformly over graphs \(G\). We apply this with \(H=K_3\).

\subsubsection{Theorem 2.2}\label{theorem-2.2-list-edge-coloring-proved-in-appendix-d}

Let \(B\) be a bipartite graph with maximum degree \(\Delta(B)\). If
every edge \(e\) is assigned a list \(L(e)\) with

\[
|L(e)|\ge\Delta(B),
\]

then \(B\) has a proper edge coloring choosing the color of each edge
from its list.

This is the maximum-degree case of Galvin's theorem {[}12{]}. A
self-contained proof (König coloring, stable pairings, and the kernel
lemma) is given in Appendix D, so this is not an external dependency of
the paper. The local refinement \(|L(xy)|\ge\max\{d(x),d(y)\}\) of
Borodin, Kostochka, and Woodall {[}2{]} is not needed: hypothesis (7.2)
already bounds every list by the maximum degree of the gain graph.

\subsubsection{Theorem 2.3}\label{theorem-2.3-dense-triangle-decomposition}

For every \(\varepsilon>0\), every sufficiently large triangle-divisible
graph \(H\) with

\[
\delta(H)\ge(0.9+\varepsilon)|V(H)|
\]

admits a triangle decomposition. This follows from Dross's fractional
triangle decomposition theorem (minimum degree \(0.9v\) suffices
fractionally) combined with the iterative-absorption theorem of Barber,
Kühn, Lo, and Osthus, which converts the fractional threshold into an
exact decomposition threshold for triangle-divisible graphs, up to
\(\varepsilon\) and for sufficiently large order.

\begin{center}\rule{0.5\linewidth}{0.5pt}\end{center}

\section{3. The Common-Profile Linear
Program}\label{the-common-profile-linear-program}

For integers \(p,q,d\), let \(H(p,q,d)\) be the split graph with clique
\(K\), \(|K|=p\), independent set \(I\), \(|I|=q\), and a fixed set
\(N\subseteq K\), \(|N|=d\), such that every vertex of \(I\) has
neighborhood \(N\). Put \(R=K\setminus N\) and \(r=p-d\).

\subsection{3.1 Symmetric cover
variables}\label{symmetric-cover-variables}

Averaging over permutations of \(N\), \(R\), and \(I\), an optimal
fractional triangle cover may be assumed constant on the four edge
classes

\[
E(N),
\qquad E(N,I),
\qquad E(N,R),
\qquad E(R).
\]

Let their weights be \(a,b,c,e\), respectively. The triangle constraints
are

\[
a+2b\ge1,
\qquad
3a\ge1,
\tag{3.1}
\]

\[
a+2c\ge1,
\qquad
2c+e\ge1,
\qquad
3e\ge1.
\tag{3.2}
\]

Constraints corresponding to empty triangle types may be omitted; the
formula below remains valid for \(p\ge3\).

The objective is

\[
\binom d2a+qdb+dr\,c+\binom r2e.
\tag{3.3}
\]

\subsection{3.2 Exact solution}\label{exact-solution}

\subsubsection{Theorem 3.1}\label{theorem-3.1-common-profile-formula}

For \(p\ge3\),

\[
\boxed{
\nu_3^*(H(p,q,d))
=
F(p,q,d),
}
\tag{3.4}
\]

where

\[
\boxed{
F(p,q,d)
=
\min\left\{
\frac{\binom p2+qd}{3},
\binom d2+\binom r2,
\binom d2+\frac{dr+\binom r2}{3}
\right\}.
}
\tag{3.5}
\]

\subsubsection{Proof}\label{proof}

At an optimum,

\[
b=\frac{1-a}{2},
\qquad
c=\frac{1-\min\{a,e\}}2,
\qquad
\frac13\le a,e\le1.
\]

If \(a\le e\), decreasing \(e\) to \(a\) cannot increase the objective.
It therefore suffices to minimize on

\[
\frac13\le e\le a\le1.
\]

The objective is affine on this triangle, so a minimum occurs at

\[
(a,e)=\left(\frac13,\frac13\right),
\qquad
(1,1),
\qquad
\left(1,\frac13\right).
\]

The corresponding values are precisely the three expressions in (3.5).
Duality completes the proof. \(\square\)

\subsection{3.3 Interpretation of the three
covers}\label{interpretation-of-the-three-covers}

The minimum in (3.5) is easier to remember through the three vertices of
the reduced two-variable program. Each vertex represents a distinct way
to pay for the triangle constraints.

\begin{longtable}[]{@{}
  >{\raggedright\arraybackslash}p{(\columnwidth - 6\tabcolsep) * \real{0.2143}}
  >{\raggedleft\arraybackslash}p{(\columnwidth - 6\tabcolsep) * \real{0.2857}}
  >{\raggedleft\arraybackslash}p{(\columnwidth - 6\tabcolsep) * \real{0.2857}}
  >{\raggedright\arraybackslash}p{(\columnwidth - 6\tabcolsep) * \real{0.2143}}@{}}
\toprule\noalign{}
\begin{minipage}[b]{\linewidth}\raggedright
Cover pattern
\end{minipage} & \begin{minipage}[b]{\linewidth}\raggedleft
Reduced vertex \((a,e)\)
\end{minipage} & \begin{minipage}[b]{\linewidth}\raggedleft
Value
\end{minipage} & \begin{minipage}[b]{\linewidth}\raggedright
Geometric description
\end{minipage} \\
\midrule\noalign{}
\endhead
\bottomrule\noalign{}
\endlastfoot
\textbf{Uniform} & \((1/3,1/3)\) & \(\bigl(\binom p2+qd\bigr)/3\) & All
relevant edge orbits are covered at the uniform fractional rate. \\
\textbf{Separated} & \((1,1)\) & \(\binom d2+\binom r2\) & The two
clique blocks \(N\) and \(R\) are paid internally, with no weight on the
crossing orbit. \\
\textbf{Hot neighborhood} & \((1,1/3)\) &
\(\binom d2+\bigl(dr+\binom r2\bigr)/3\) & The common neighborhood \(N\)
is paid in full, while the residual clique structure remains
fractional. \\
\end{longtable}

The labels are descriptive rather than additional definitions. In later
calculations we refer to the corresponding first, second, and third
branches of \(F\). This finite trichotomy is the source of the global
fractional margin.

\begin{center}\rule{0.5\linewidth}{0.5pt}\end{center}

\section{4. Fractional Cloning and the Exact Fractional
Margin}\label{fractional-cloning-and-the-exact-fractional-margin}

\subsection{4.1 Fractional cloning}\label{fractional-cloning}

The cloning operation here is the cover-side analogue of the clone
symmetrization used in Paper II. There, vertices or clone classes are
copied pairwise to simplify a chordal graph. Here, one independent-set
profile is copied \(q\) times in a single averaging step; the graph
\(H(p,q,d_i)\) records the resulting common profile.

\subsubsection{Lemma 4.1}\label{lemma-4.1-fractional-cloning-bound}

For every split graph with \(q\ge1\),

\[
\boxed{
\nu_3^*(G)
\ge
\frac1q\sum_{i=1}^qF(p,q,d_i).
}
\tag{4.1}
\]

\subsubsection{Proof}\label{proof-1}

Let \(y\) be any fractional triangle cover of \(G\). Put

\[
A=\sum_{e\in E(K)}y_e
\]

and

\[
B_i=\sum_{x\in N_i}y_{v_ix}.
\]

Replace \(v_i\) by \(q\) independent clones, all with neighborhood
\(N_i\), and give each clone the incident weights of \(v_i\). Together
with the original clique-edge weights, this is a fractional cover of
\(H(p,q,d_i)\) of weight \(A+qB_i\). Hence

\[
A+qB_i\ge F(p,q,d_i).
\]

Summing over \(i\) gives

\[
q\left(A+\sum_iB_i\right)
\ge
\sum_iF(p,q,d_i).
\]

Minimizing over covers proves the lemma. \(\square\)

\subsection{4.2 The exact margin}\label{the-exact-margin}

Assume \(0<q\le2p\), set

\[
\alpha=\frac qp,
\]

and define the packing threshold

\[
T(G)=\frac12\left(|E(G)|-\frac{(p+q)^2}{6}\right).
\tag{4.2}
\]

Define

\[
\mu(\alpha)
=
\begin{cases}
\alpha^2/12,&0\le\alpha\le2/3,\\
(2-\alpha)^2/48,&2/3\le\alpha\le2.
\end{cases}
\tag{4.3}
\]

\subsubsection{Theorem 4.2}\label{theorem-4.2-unified-fractional-margin}

\[
\boxed{
\nu_3^*(G)
\ge
T(G)+\mu(\alpha)p^2-\frac p4.
}
\tag{4.4}
\]

\subsubsection{Proof}\label{proof-2}

Let

\[
C_\alpha=\frac{2-2\alpha-\alpha^2}{12}.
\]

For each branch of \(F(p,q,d)\), completion of squares gives

\[
F(p,q,d)
\ge
\frac{qd}{2}+C_\alpha p^2+\mu(\alpha)p^2-\frac p2.
\tag{4.5}
\]

The three normalized residual minima are

\[
\frac{\alpha^2}{12},
\qquad
\frac{(2-\alpha)^2}{48},
\]

and

\[
\begin{cases}
\alpha(8-5\alpha)/48,&\alpha\le4/3,\\
(2-\alpha)^2/12,&\alpha\ge4/3.
\end{cases}
\]

The third is never below the minimum of the first two. Averaging (4.5)
through Lemma 4.1 gives

\[
\nu_3^*(G)
\ge
\frac12\sum_i d_i+C_\alpha p^2+\mu(\alpha)p^2-\frac p2.
\]

Since

\[
T(G)=\frac12\sum_i d_i+C_\alpha p^2-\frac p4,
\]

we obtain (4.4). \(\square\)

\subsection{4.3 Bulk consequence}\label{bulk-consequence}

If

\[
\varepsilon\le\alpha\le2-\varepsilon,
\]

then \(\mu(\alpha)\ge c_\varepsilon>0\). Theorem 4.2 and Haxell--Rödl
imply

\[
\nu_3(G)\ge T(G)
\]

for all sufficiently large graphs in this regime. Hence

\[
|E(G)|-2\nu_3(G)
\le
\frac{n^2}{6}.
\tag{4.6}
\]

\begin{center}\rule{0.5\linewidth}{0.5pt}\end{center}

\section{\texorpdfstring{5. Factorization Rounding Near
\(q=2p\)}{5. Factorization Rounding Near q=2p}}\label{factorization-rounding-near-q2p}

Let

\[
r_p=\chi'(K_p).
\]

\subsection{5.1 One-factor averaging}\label{one-factor-averaging}

\subsubsection{Lemma 5.1}\label{lemma-5.1}

If \(q\ge r_p\), then

\[
\boxed{
\nu_3(G)
\ge
\frac1q\sum_i\binom{d_i}{2}.
}
\tag{5.1}
\]

\subsubsection{Proof}\label{proof-3}

Factor \(K_p\) into \(r_p\) matchings. Assign the factors injectively
and uniformly to vertices of \(I\). In the factor assigned to \(v_i\),
retain only edges with both endpoints in \(N_i\). The retained edges
form valid edge-disjoint \(KKI\) triangles.

The expected number retained is

\[
\frac1q\sum_i\sum_{j=1}^{r_p}|F_j\cap E(K[N_i])|
=
\frac1q\sum_i\binom{d_i}{2}.
\]

Some assignment attains at least the expectation. \(\square\)

Writing \(q=2p-s\), substitution gives

\[
\Phi(G)
\le
\frac{n^2}{6}+\frac p2-\frac{s^2}{6}
+
\frac{(s-1)M-S_2}{q}.
\tag{5.2}
\]

Using \(S_2\ge M^2/q\) and maximizing the resulting parabola yields

\[
\boxed{
\Phi(G)
\le
\frac{n^2}{6}+\frac p2+\frac{s^2-6s+3}{12}.
}
\tag{5.3}
\]

Thus \(s=O(\sqrt p)\) is closed with linear error.

\subsection{5.2 Double-factor rounding}\label{double-factor-rounding}

For a clique edge \(e=xy\), let

\[
b_e=|\{i:\{x,y\}\not\subseteq N_i\}|.
\]

Put

\[
h=\min\{r_p,q-r_p\},
\qquad
\delta=\frac h{r_p}.
\]

Choose the set of \(h\) factors that receive two distinct independent
vertices uniformly at random, give one independent vertex to every other
factor, and assign the independent vertices to the resulting positions
uniformly and injectively. All expectations below are taken over both
random choices; some deterministic choice attains at least the expected
gain.

\subsubsection{Lemma 5.2}\label{lemma-5.2-double-factor-inequality}

If \(q\ge r_p\), then

\[
\boxed{
\begin{aligned}
\Phi(G)
\le{}&
\frac{n^2}{6}+\frac p2-\frac{s^2}{6}
+
\frac{(s-1)M-S_2}{q}\\
&-
\frac{2\delta V}{q(q-1)},
\end{aligned}
}
\tag{5.4}
\]

where

\[
V=\sum_{e\in E(K)}b_e(q-b_e).
\tag{5.5}
\]

\subsubsection{Proof}\label{proof-4}

If the factor containing \(e\) receives one vertex, \(e\) is lost with
probability \(b_e/q\). If it receives two vertices, it is lost only when
both are bad, with probability

\[
\frac{b_e(b_e-1)}{q(q-1)}.
\]

Thus the expected number \(U\) of lost clique edges is

\[
U
=
\frac1q\sum_e b_e
-
\frac{\delta}{q(q-1)}\sum_e b_e(q-b_e).
\]

Because every factor is a matching, each retained edge may be assigned
to one of its admitting vertices without creating an edge collision.
Substitution gives (5.4). \(\square\)

\begin{center}\rule{0.5\linewidth}{0.5pt}\end{center}

\section{6. Quantitative Polarization}\label{quantitative-polarization}

For every \(i\), define

\[
\mathcal B_i=\{e\in E(K):e\cap S_i\ne\varnothing\}.
\]

Then

\[
V=\sum_{i,j}|\mathcal B_i\setminus\mathcal B_j|.
\tag{6.1}
\]

Let

\[
m=\max_i|S_i|.
\]

\subsubsection{Lemma 6.1}\label{lemma-6.1-polarization-inequality}

If \(2p-3m-1\ge0\), then

\[
\boxed{
V
\ge
\frac{2p-3m-1}{4}
\sum_{i,j}|S_i\triangle S_j|.
}
\tag{6.2}
\]

\subsubsection{Proof}\label{proof-5}

Put

\[
a_{ij}=|S_i\setminus S_j|.
\]

The edges in \(\mathcal B_i\setminus\mathcal B_j\) are precisely the
edges of \(K_p-S_j\) meeting \(S_i\setminus S_j\). Therefore

\[
|\mathcal B_i\setminus\mathcal B_j|
=
\frac{a_{ij}\bigl(2(p-|S_j|)-a_{ij}-1\bigr)}2.
\]

Since \(|S_j|\le m\) and \(a_{ij}\le m\),

\[
|\mathcal B_i\setminus\mathcal B_j|
\ge
\frac{2p-3m-1}{2}|S_i\setminus S_j|.
\]

Sum over ordered pairs and use

\[
\sum_{i,j}|S_i\setminus S_j|
=
\frac12\sum_{i,j}|S_i\triangle S_j|.
\]

This proves (6.2). \(\square\)

\begin{center}\rule{0.5\linewidth}{0.5pt}\end{center}

\section{7. Shifted-Center Gain
Completion}\label{shifted-center-gain-completion}

Fix \(R\subseteq K\). Put

\[
\rho=|R|,
\qquad
Q=K\setminus R,
\qquad
b=|Q|.
\]

For each \(i\), define

\[
T_i=S_i\setminus R,
\qquad
t_i=|T_i|,
\]

and

\[
G_i=R\setminus S_i,
\qquad
g_i=|G_i|.
\]

Set

\[
A_R=\sum_i t_i,
\qquad
A_{2,R}=\sum_i t_i^2,
\qquad
B_R=\sum_i g_i.
\]

Let

\[
r_b=\chi'(K_b),
\qquad
u=q-r_b.
\]

Assume

\[
b\ge2,
\qquad
q\ge r_b,
\qquad
b\ge\chi'(K_\rho),
\tag{7.1}
\]

and

\[
b-t_i\ge\max\{\rho,u\}
\qquad
\text{for all }i.
\tag{7.2}
\]

Define

\[
\theta_R=\frac{\max\{\rho-1,0\}}{b}
\]

and

\[
\kappa_R
=
1-2(1-\theta_R)\frac{u}{q}.
\tag{7.3}
\]

\subsection{\texorpdfstring{7.1 Packing inside
\(Q\)}{7.1 Packing inside Q}}\label{packing-inside-q}

Reserve a set \(U\subseteq I\) of \(u\) vertices; note that
\(|I\setminus U|=q-u=r_b\) exactly. Assign the remaining \(r_b\)
vertices bijectively to a factorization of \(K[Q]\).

For \(v_i\), the number of unavailable edges of \(K[Q]\) is

\[
\beta_i
=
\binom b2-\binom{b-t_i}{2}.
\tag{7.4}
\]

For fixed \(U\), averaging over assignments gives at least

\[
\binom b2-\frac1{r_b}\sum_{i\notin U}\beta_i
\tag{7.5}
\]

edge-disjoint \(QQI\) triangles.

\subsection{7.2 The gain graph}\label{the-gain-graph}

Construct a bipartite graph with parts \(U\) and \(R\), joining \(v_i\)
to \(r\) when \(r\in G_i\). Assign the list

\[
L(v_ir)=N_i\cap Q.
\]

By (7.2), every list satisfies

\[
|L(v_ir)|=b-t_i
\ge
\max\{\rho,u\}
\ge
\Delta,
\]

where \(\Delta\) is the maximum degree of the gain graph, since
\(d(v_i)=g_i\le\rho\) and \(d(r)\le|U|=u\). Theorem 2.2, proved
self-contained in Appendix D, gives a proper list edge coloring. If
\(v_ir\) receives color \(z\in Q\), take the triangle \(v_irz\). This
gives

\[
B_U=\sum_{i\in U}g_i
\]

edge-disjoint \(IRQ\) triangles.

\subsection{\texorpdfstring{7.3 Completion inside
\(R\)}{7.3 Completion inside R}}\label{completion-inside-r}

For each \(z\in Q\), let \(U_z\subseteq R\) be the vertices \(r\) for
which \(rz\) was used by an \(IRQ\) triangle.

If \(\rho\le1\), the graph \(K[R]\) has no edges, so this completion
contributes no \(RRQ\) triangles. Suppose therefore that \(\rho\ge2\).
Factor \(K[R]\), inject its factors into the colors \(z\in Q\), and
delete from the factor assigned to \(z\) every edge incident with
\(U_z\). Averaging over the injections loses at most

\[
\frac{\rho-1}{b}\sum_z|U_z|
=
\theta_RB_U
\]

edges. Hence, in all cases, we obtain at least

\[
\binom\rho2-\theta_RB_U
\]

additional \(RRQ\) triangles.

The three families \(QQI\), \(IRQ\), and \(RRQ\) are edge-disjoint. In
particular, the forbidden-color deletion prevents an \(IRQ\) triangle
and an \(RRQ\) triangle from sharing an edge \(rz\).

\subsection{7.4 The centered inequality}\label{the-centered-inequality}

Choose \(U\), \(|U|=u\), maximizing

\[
\sum_{i\in U}
\left(
\frac{2\beta_i}{r_b}+2(1-\theta_R)g_i
\right).
\]

The best \(u\) terms have sum at least \(u/q\) times the total. Using

\[
2\sum_i\beta_i
=(2b-1)A_R-A_{2,R},
\]

we obtain the principal local lemma.

\subsubsection{Lemma 7.1}\label{lemma-7.1-reserved-gain-shifted-center-inequality}

Under (7.1)--(7.2),

\[
\boxed{
\begin{aligned}
\Phi(G)
\le{}&
\frac{n^2}{6}+\frac p2-\frac{s^2}{6}
+s\rho-2\rho^2\\
&+
\kappa_RB_R
+
\frac{(s-2\rho-1)A_R-A_{2,R}}q.
\end{aligned}
}
\tag{7.6}
\]

\begin{center}\rule{0.5\linewidth}{0.5pt}\end{center}

\section{8. The Sparse-Independent
Regime}\label{the-sparse-independent-regime}

Assume \(q=o(p)\). We prove the required bound directly and integrally.

\subsection{8.1 Minimal-counterexample degree
bound}\label{minimal-counterexample-degree-bound}

In the global contradiction argument, a minimal counterexample with
penalty \(kn\) satisfies, for every \(v\in I\),

\[
d(v)>
\frac{2n-1}{6}+k.
\tag{8.1}
\]

Since \(q=o(p)\), this implies eventually

\[
d(v)\ge2q+2.
\tag{8.2}
\]

\subsection{8.2 Successive matchings}\label{successive-matchings}

Order \(I=\{v_1,\ldots,v_q\}\). Choose successively edge-disjoint
matchings

\[
M_i\subseteq K[N_i]
\]

of size \(\lfloor d_i/2\rfloor\).

Before choosing \(M_i\), each vertex has lost at most \(i-1\) incident
edges, so the available graph on \(N_i\) has minimum degree at least

\[
d_i-i\ge d_i/2.
\]

Dirac's theorem supplies a Hamilton cycle and hence a matching of the
required size.

Let

\[
F=\bigcup_iM_i.
\]

Then

\[
|F|
\ge
\frac12\sum_i d_i-\frac q2,
\qquad
\Delta(F)\le q.
\tag{8.3}
\]

The edges of \(M_i\) yield \(|M_i|\) edge-disjoint \(KKI\) triangles
with apex \(v_i\).

\subsection{8.3 Divisibility correction}\label{divisibility-correction}

The residual clique graph

\[
R_0=K_p-F
\]

satisfies

\[
\delta(R_0)
\ge
p-1-q
=(1-o(1))p.
\tag{8.4}
\]

In particular, for all sufficiently large members of the sequence,
\(\delta(R_0)>p/2\). Dirac's theorem therefore gives a Hamilton cycle,
and hence a Hamilton path

\[
P=x_1x_2\cdots x_p
\]

contained in \(R_0\). Let \(O\) be the set of odd-degree vertices of
\(R_0\). By the handshaking lemma, \(|O|\) is even. Define the path
subgraph \(J\subseteq P\) by

\[
x_jx_{j+1}\in E(J)
\quad\Longleftrightarrow\quad
|O\cap\{x_1,\ldots,x_j\}|\text{ is odd}.
\tag{8.5}
\]

For an internal vertex \(x_j\), the parity of its degree in \(J\) is the
change in prefix parity between positions \(j-1\) and \(j\), and is
therefore \(1\) precisely when \(x_j\in O\). The same conclusion holds
at the two endpoints because \(|O|\) is even. Consequently,

\[
\operatorname{Odd}(J)=O,
\qquad
|E(J)|\le p-1,
\qquad
\Delta(J)\le2.
\tag{8.6}
\]

Thus

\[
R_1:=R_0-J
\]

has all degrees even, and

\[
\delta(R_1)\ge p-1-q-2.
\tag{8.7}
\]

Since \(q=o(p)\), eventually \(\delta(R_1)>3p/4\). By Turán's theorem,
\(R_1\) contains a copy of \(K_5\). Choose a \(C_4\) and a \(C_5\)
inside this fixed \(K_5\). According as

\[
|E(R_1)|\equiv0,1,2\pmod3,
\]

remove respectively nothing, the \(C_4\), or the \(C_5\). Let \(C\)
denote the removed cycle, allowing \(C=\varnothing\), and put

\[
H:=R_1-E(C).
\]

Every vertex loses either zero or two incident edges when a cycle is
removed, so all degrees of \(H\) remain even. Moreover,
\(|E(C)|\equiv |E(R_1)|\pmod3\), and hence

\[
|E(H)|\equiv0\pmod3.
\]

Therefore \(H\) is triangle-divisible. The total loss of degree in
passing from \(R_0\) to \(H\) is at most four, so

\[
\delta(H)\ge p-1-q-4.
\tag{8.8}
\]

To apply Theorem 2.3 with a fixed parameter, take for example
\(\varepsilon_0=1/100\). Since \(q=o(p)\), eventually \(q\le p/20\), and
then, for sufficiently large \(p\),

\[
\delta(H)\ge p-1-q-4\ge0.91p=(0.9+\varepsilon_0)|V(H)|.
\tag{8.9}
\]

Theorem 2.3 now gives an exact triangle decomposition of \(H\). Notice
that the decomposition theorem is applied on the original \(p\)-vertex
set: no vertices are deleted during the correction. The entire
correction removes at most \(p+4\) edges.

\subsection{8.4 Packing size}\label{packing-size}

The combined packing has size

\[
\begin{aligned}
\nu_3(G)
&\ge
|F|+\frac{\binom p2-|F|-O(p)}3\\
&=
\frac13\binom p2+\frac23|F|-O(p)\\
&\ge
\frac13\binom p2+\frac13\sum_i d_i-O(p+q).
\end{aligned}
\tag{8.10}
\]

Consequently,

\[
\begin{aligned}
\Phi(G)
&\le
\frac13\binom p2+\frac13\sum_i d_i+O(p+q)\\
&\le
\frac13\binom p2+\frac{pq}{3}+O(p+q)\\
&=
\frac{(p+q)^2}{6}-\frac{p+q^2}{6}+O(p+q).
\end{aligned}
\tag{8.11}
\]

This is at most \(n^2/6+O(n)\).

\begin{center}\rule{0.5\linewidth}{0.5pt}\end{center}

\section{9. Proof of the Main Theorem}\label{proof-of-the-main-theorem}

We now prove Theorem 1.1 by contradiction.

Suppose no absolute constant works. For every positive integer \(k\),
choose a split graph \(G_k\), minimal in its number \(n_k\) of vertices,
such that

\[
\Phi(G_k)
>
\frac{n_k^2}{6}+kn_k.
\tag{9.1}
\]

Necessarily \(n_k\to\infty\). For every \(v\in I(G_k)\), minimality
gives

\[
\Phi(G_k)
\le
\Phi(G_k-v)+d(v),
\]

and therefore

\[
\boxed{
d(v)>
\frac{2n_k-1}{6}+k.
}
\tag{9.2}
\]

First, if \(q_k\ge2p_k-1\), Lemma 5.1 gives

\[
\Phi(G_k)
\le
\frac{n_k^2}{6}+\frac{p_k}{2},
\]

contradicting (9.1) for large \(k\). We may therefore assume
\(q_k<2p_k-1\) and pass to a subsequence according to

\[
\alpha_k=\frac{q_k}{p_k}\in[0,2).
\]

\subsection{9.1 Bulk regime}\label{bulk-regime}

Suppose that for some \(\varepsilon>0\),

\[
\varepsilon\le\alpha_k\le2-\varepsilon.
\]

Theorem 4.2 gives

\[
\nu_3^*(G_k)
\ge
T(G_k)+c_\varepsilon p_k^2-O(p_k).
\]

Haxell--Rödl gives

\[
\nu_3(G_k)
\ge
\nu_3^*(G_k)-o(n_k^2).
\]

For sufficiently large \(k\), the quadratic margin dominates the
integrality loss, so \(\nu_3(G_k)\ge T(G_k)\), contrary to (9.1).

\subsection{\texorpdfstring{9.2 The endpoint
\(\alpha_k\to0\)}{9.2 The endpoint \textbackslash alpha\_k\textbackslash to0}}\label{the-endpoint-alpha_kto0}

This is closed by Section 8, which gives an absolute \(C_0\) such that

\[
\Phi(G_k)
\le
\frac{n_k^2}{6}+C_0n_k.
\]

For \(k>C_0\), this contradicts (9.1).

\subsection{\texorpdfstring{9.3 The endpoint
\(\alpha_k\to2\)}{9.3 The endpoint \textbackslash alpha\_k\textbackslash to2}}\label{the-endpoint-alpha_kto2}

Write

\[
q=2p-s,
\qquad
s=o(p).
\]

If \(s=O(\sqrt p)\), inequality (5.3) gives the required linear-error
bound. Assume henceforth

\[
\sqrt p\ll s=o(p).
\tag{9.3}
\]

For all sufficiently large members of the sequence,

\[
p\ge2304,
\qquad
6\sqrt p\le s\le\frac p8,
\qquad
k\ge1.
\tag{9.4}
\]

By (9.2),

\[
m_i=|S_i|
<
\frac s3+\frac16-k
\le
\frac s3-\frac56.
\]

Since every \(m_i\) is an integer, with \(m=\max_i m_i\),

\[
\boxed{3m\le s-3.}
\tag{9.5}
\]

For \(x\in K\), put

\[
a_x=|\{i:x\in S_i\}|
\]

and

\[
D=\sum_{x\in K}a_x(q-a_x).
\tag{9.6}
\]

Then

\[
\sum_{i,j}|S_i\triangle S_j|=2D.
\tag{9.7}
\]

\subsubsection{High dispersion}\label{high-dispersion}

Suppose

\[
D\ge\frac{qs^2}{12}.
\tag{9.8}
\]

By (9.5),

\[
2p-3m-1\ge2p-s+2=q+2.
\]

Lemma 6.1 and (9.7) therefore give

\[
V
\ge
\frac{q+2}{2}D
\ge
\frac q2D.
\tag{9.9}
\]

The double-factor coefficient in Lemma 5.2 satisfies

\[
\boxed{\delta\ge\frac78.}
\tag{9.10}
\]

Indeed, if \(p\) is odd then \(r_p=p\) and

\[
\delta=\frac{p-s}{p}\ge\frac78,
\]

while if \(p\) is even then \(r_p=p-1\) and

\[
\delta=\frac{p+1-s}{p-1}\ge\frac78.
\]

Moreover, \(S_2\ge M^2/q\), and (9.5) implies \(M/q\le m\le(s-3)/3\).
Since the resulting parabola is increasing on this interval,

\[
\frac{(s-1)M-S_2}{q}
\le
\frac{2s(s-3)}9.
\tag{9.11}
\]

Substituting (9.8)--(9.11) into Lemma 5.2 yields

\[
\Phi(G)-\frac{n^2}{6}
\le
\frac p2-\frac{5s^2}{288}-\frac{2s}{3}.
\tag{9.12}
\]

Since \(s^2\ge36p\),

\[
\Phi(G)-\frac{n^2}{6}
\le
-\frac p8-\frac{2s}{3}<0,
\]

a contradiction.

\subsubsection{Low dispersion}\label{low-dispersion}

Suppose

\[
D<\frac{qs^2}{12}.
\tag{9.13}
\]

By (9.7), some \(j\) satisfies

\[
\sum_i|S_i\triangle S_j|
<
\frac{s^2}{6}.
\tag{9.14}
\]

Set

\[
R=S_j,
\qquad
\rho=|R|.
\]

Then

\[
A_R+B_R<\frac{s^2}{6},
\qquad
\rho\le m\le\frac{s-3}{3}.
\tag{9.15}
\]

We verify the hypotheses of Lemma 7.1 explicitly. Since \(s\le p/8\),

\[
q=2p-s\ge\frac{15p}{8},
\qquad
b=p-\rho\ge p-\frac s3.
\]

Thus \(b\ge2\), \(q\ge r_b\), and \(b\ge\chi'(K_\rho)\). Also
\(r_b\ge b-1\), so

\[
u=q-r_b\le p-s+\rho+1.
\]

For every \(i\),

\[
2\rho+t_i+1
\le
3m+1
\le
s-2.
\]

Consequently,

\[
b-t_i\ge\max\{\rho,u\},
\]

and Lemma 7.1 applies.

If \(\rho=0\), then \(B_R=0\). If \(\rho\ge1\), then

\[
\theta_R
\le
\frac{\rho}{p-\rho}
\le
\frac{8s}{23p}.
\]

Using \(u=q-r_b\), \(r_b\le b\), and \(q\ge15p/8\), we obtain

\[
\kappa_R
=1-2(1-\theta_R)\frac uq
\le
\frac sq+2\theta_R
\le
\frac{5s}{4p}.
\tag{9.16}
\]

Likewise,

\[
\frac{(s-2\rho-1)_+}{q}
\le
\frac{5s}{4p}.
\tag{9.17}
\]

By (9.15), the total positive deviation in (7.6) is therefore at most

\[
\frac{5s}{4p}(A_R+B_R)
<
\frac{5s^3}{24p}
\le
\frac{5s^2}{192}.
\tag{9.18}
\]

Meanwhile,

\[
\frac{s^2}{6}-s\rho+2\rho^2
=
2\left(\rho-\frac s4\right)^2+\frac{s^2}{24}
\ge
\frac{s^2}{24}.
\tag{9.19}
\]

Lemma 7.1 now gives

\[
\Phi(G)-\frac{n^2}{6}
\le
\frac p2-\frac{s^2}{64}.
\tag{9.20}
\]

Since \(s^2\ge36p\), the right-hand side is at most \(-p/16\), again a
contradiction.

All possible subsequences are impossible. Theorem 1.1 follows.
\(\square\)

\begin{center}\rule{0.5\linewidth}{0.5pt}\end{center}

\section{10. Corollaries}\label{corollaries}

\subsection{10.1 Clique partitions}\label{clique-partitions}

Corollary 1.2 follows immediately from

\[
\operatorname{cp}(G)
\le
|E(G)|-2\nu_3(G).
\]

\subsection{10.2 Sharpness of the leading
term}\label{sharpness-of-the-leading-term}

The family

\[
K_p\vee\overline K_{2p}
\]

has \(n=3p\). A factorization of \(K_p\) assigned to independent
vertices packs every clique edge into a \(KKI\) triangle. The resulting
triangle-packing expression is

\[
|E|-2\nu_3
=
\frac{n^2}{6}+\frac n6.
\tag{10.1}
\]

The Erdős--Ordman--Zalcstein clique-partition construction on this
family shows that the coefficient \(1/6\) cannot be improved.

\subsection{10.3 Localization of difficult
sequences}\label{localization-of-difficult-sequences}

Proposition 1.3 records the localization statement supplied by the
proof. It should be understood as a structural description of a
hypothetical failure of every linear-error estimate, rather than as an
additional extremal theorem after Theorem 1.1. The bulk margin excludes
ratios bounded away from \(0\) and \(2\), the sparse construction
excludes \(|I|/|K|\to0\), and averaged factorization excludes the short
corridor near \(|I|=2|K|\). The only remaining location is the
mesoscopic corridor

\[
|I|=2|K|-s,
\qquad
\sqrt{|K|}\ll s=o(|K|),
\]

which is then eliminated by the dispersion dichotomy.

\subsection{10.4 Exact common-profile fractional
value}\label{exact-common-profile-fractional-value}

Theorem 3.1 is independently useful: it gives the exact fractional
triangle packing for every split graph in which all independent vertices
have one common neighborhood, including the regime with too few
independent vertices to color all neighborhood edges integrally. We
state it as a stand-alone corollary.

\subsubsection{Corollary 10.4}\label{corollary-10.4-exact-common-neighborhood-triangle-packing}

For all integers \(p\ge3\), \(q\ge0\) and \(0\le d\le p\), the graph
\(H(p,q,d)\) (clique of order \(p\); \(q\) independent vertices, each
with the same neighborhood of size \(d\)) satisfies \[
\nu_3^*(H(p,q,d))=F(p,q,d)=\min\left\{\tfrac{\binom p2+qd}{3},\ \binom d2+\binom{p-d}2,\ \binom d2+\tfrac{d(p-d)+\binom{p-d}2}{3}\right\}.
\] This is an exact closed form for the fractional triangle-packing
number of an entire one-parameter family, valid for every \(d\)
including the small-\(q\) regime where no integral factorization covers
the neighborhood; it may be of independent interest outside the
clique-partition problem.

\subsubsection{Corollary 10.4b}\label{corollary-10.4b-threshold-graphs}

Since \(K_p\vee\overline K_{2p}\) is a threshold graph and every
threshold graph is split, Theorem 1.1 gives, in particular, the
linear-error bound \(\operatorname{cp}(G)\le n^2/6+Cn\) for all
threshold graphs, with the same sharp leading constant witnessed inside
that subclass.

\subsection{\texorpdfstring{10.5 Localized effectivity near
\(q=2p\)}{10.5 Localized effectivity near q=2p}}\label{localized-effectivity-near-q2p}

\subsubsection{Proposition 10.1}\label{proposition-10.1-explicit-corridor-bounds}

Let \(q=2p-s\) and \(n=p+q\).

\begin{enumerate}
\def\labelenumi{\arabic{enumi}.}
\item
  If \[
  p\ge36,
  \qquad
  0\le s\le6\sqrt p,
  \] then \[
  \boxed{
  \Phi(G)\le\frac{n^2}{6}+2n.
  }
  \tag{10.2}
  \]
\item
  If \[
  p\ge2304,
  \qquad
  6\sqrt p\le s\le\frac p8,
  \] and every \(v\in I\) satisfies \[
  d(v)>\frac{2n-1}{6}+1,
  \] then \[
  \boxed{
  \Phi(G)\le\frac{n^2}{6}.
  }
  \tag{10.3}
  \]
\end{enumerate}

Consequently, for every \(C\ge2\), a minimum-order counterexample to

\[
\Phi(G)\le\frac{n^2}{6}+Cn
\]

cannot satisfy

\[
p\ge2304,
\qquad
0\le s\le\frac p8.
\]

\subsubsection{Proof}\label{proof-6}

For the short corridor, inequality (5.3) and \(s^2\le36p\) give

\[
\Phi(G)-\frac{n^2}{6}
\le
\frac{7p}{2}+\frac14
\le
2n,
\]

because \(p\ge36\) implies \(s\le p\) and hence \(n=3p-s\ge2p\).

The mesoscopic assertion is exactly the quantitative high/low-dispersion
calculation in Section 9.3. The final consequence follows from the
degree inequality obtained by deleting an independent vertex from a
minimum-order counterexample. \(\square\)

\begin{center}\rule{0.5\linewidth}{0.5pt}\end{center}

\section{11. Discussion}\label{discussion}

\subsection{11.1 From fractional asymptotics to a linear-error integral
theorem}\label{from-fractional-asymptotics-to-a-linear-error-integral-theorem}

The central difficulty is that the general theorem

\[
\nu_3^*(G)-\nu_3(G)=o(n^2)
\]

has no linear rate. The present proof avoids demanding one uniform
rounding mechanism.

\begin{itemize}
\tightlist
\item
  In the bulk, the common-profile LP produces quadratic slack.
\item
  Near \(q=2p\), explicit factorization structures replace general
  rounding.
\item
  Near \(q=0\), the residual clique becomes dense enough for exact
  decomposition.
\end{itemize}

This regime-dependent strategy is weaker than a universal linear
integrality-gap theorem but sufficient for the extremal clique-partition
problem.

\subsection{11.2 Why residual maximum degree is not the right
invariant}\label{why-residual-maximum-degree-is-not-the-right-invariant}

Early exploratory work attempted to prove that an optimal packing leaves
a clique residual of bounded maximum degree. Numerical behavior depended
strongly on the packing algorithm, and no such property is needed in the
final proof.

The effective invariants are instead:

\begin{itemize}
\tightlist
\item
  the common-profile fractional margin;
\item
  the variance term \(S_2\);
\item
  the polarization energy \(V\);
\item
  the centered deviations \(A_R,B_R\);
\item
  divisibility of the dense residual.
\end{itemize}

\subsection{11.3 Localized effectivity and remaining
non-effectivity}\label{localized-effectivity-and-remaining-non-effectivity}

Proposition 10.1 shows that the entire near-extremal corridor treated by
factorization, polarization, and shifted-center completion is
quantitatively effective. One may take

\[
C_{\mathrm{corr}}=2,
\qquad
p_0=2304,
\qquad
s_0(p)=6\sqrt p,
\qquad
\eta_0=\frac18.
\]

The global constant \(C\) remains non-effective for two independent
reasons.

\begin{enumerate}
\def\labelenumi{\arabic{enumi}.}
\tightlist
\item
  The Haxell--Rödl theorem is used through an asymptotic \(o(n^2)\)
  statement in the bulk.
\item
  Dense triangle decomposition theorems are invoked with unspecified
  sufficiently-large thresholds in the sparse-independent regime.
\end{enumerate}

Thus no hidden non-effectivity remains in Sections 5--7; it comes only
from the two external asymptotic inputs above. Moreover, after Appendix
D, Sections 5--7 are self-contained: Proposition 10.1 depends on no
external theorem whatsoever.

\subsection{11.4 Relationship with the fractional extremal
theorem}\label{relationship-with-the-fractional-extremal-theorem}

The complete-split extremal theorem determines the global fractional
maximum of

\[
|E(G)|-2\nu_3^*(G)
\]

over both split and chordal graphs. The common-profile LP in the present
paper serves a different purpose: it is a profile-sensitive local lower
bound for the replicated graph, and its quantitative slack absorbs
integrality loss away from the extremal corridor.

Thus the global fractional extremal calculation and the local
common-profile calculation are compatible but logically distinct. The
present proof is self-contained apart from the external theorems stated
in Section 2.

\textbf{Remark (cloning across the series).} Lemma 4.1 and the
clone-symmetrization step of Paper II use the same underlying operation
in different roles. Paper II uses pairwise cloning to move a chordal
graph toward a complete-split terminal graph. Lemma 4.1 uses \(q\)-fold
cloning on the cover side to compare an arbitrary profile with a
common-profile graph. The shared term \emph{cloning} is therefore
retained, while the qualifier \emph{fractional} distinguishes the
present inequality.

\subsection{11.5 Relationship with the chordal linear-error
problem}\label{relationship-with-the-chordal-linear-error-problem}

The complete-split reduction settles the asymptotically sharp chordal
bound

\[
\operatorname{cp}(G)
\le
\left(\frac16+o(1)\right)n^2.
\]

The stronger estimate

\[
\operatorname{cp}(G)
\le
\frac{n^2}{6}+O(n)
\]

remains open for general chordal graphs. The mechanisms developed here
suggest a concrete three-part program for that problem.

\begin{enumerate}
\def\labelenumi{\arabic{enumi}.}
\item
  \textbf{Stability of the fractional extremal theorem.} Upgrade the
  exact chordal fractional theorem of Paper II to a stability statement:
  every chordal graph whose value is within \(\delta n^2\) of the
  fractional maximum is within \(\varepsilon n^2\) edge edits of a
  maximizing complete split graph. Away from near-extremality, the
  quadratic fractional slack absorbs the general integrality loss
  exactly as in Section 4.
\item
  \textbf{Robust corridor mechanisms.} Extend the factorization,
  polarization, and shifted-center arguments of Sections 5--7 from split
  graphs to \(\varepsilon n^2\)-perturbations of complete split graphs,
  so that they apply to the near-extremal graphs produced by stability.
\item
  \textbf{Clique-tree assembly.} Patch the local estimates across the
  clique tree by an edge-ownership ledger that preserves separator
  capacities, in the spirit of shifted-center completion, so that no
  edge is charged twice.
\end{enumerate}

A chordal linear-error proof would still have to preserve ownership and
separator capacities across the clique tree; the ledger of step 3 is
designed for exactly that purpose.

A companion paper in preparation investigates this program. Its intended
clique-separator recursion uses the split theorem proved here as the
base case and is designed so that the chordal regimes mirror the
corresponding ingredients of Sections 3--9: common-profile evaluation,
fractional cloning, factorization, and corridor analysis. Until that
companion argument is complete, this paragraph records the proposed
interface rather than a theorem of the present paper.

\subsection{11.6 Formal-verification perimeter and current
status}\label{formal-verification-perimeter-and-current-status}

Papers I and II of the series are formally verified in Lean 4 with
Mathlib. Paper III is being verified under the same release protocol.
The Lean project for the present manuscript is built against Mathlib
v4.28.0 and currently compiles as a complete dependency graph.

The formal development separates two layers.

\begin{itemize}
\tightlist
\item
  \textbf{Layer X} contains exactly the two external asymptotic inputs,
  Theorems 2.1 and 2.3, retained as named axioms with their hypotheses
  stated explicitly.
\item
  \textbf{Layer E} contains the finite mathematical core: Theorems 3.1
  and 4.2; Lemmas 4.1--7.1; the divisibility correction of Appendix B;
  the list edge coloring development of Appendix D; the quantitative
  chains of Sections 9.3 and 10.5; Proposition 10.1; and the final
  assembly.
\end{itemize}

At the date of this draft, the project has eight localized
\texttt{sorry} obligations. They are confined to: the small degenerate
cases of the common-profile LP; the one- and double-factor rounding
block of Section 5; the middle-corridor conclusion of Proposition 10.1;
one matching helper in the self-contained Galvin/König development; the
three-family packing construction underlying Section 7.2; and the
sparse-regime packing estimate that invokes Theorem 2.3. The statement
and dependency assembly of Theorem 1.1 already compile, but the theorem
remains transitively open until these leaves are discharged.

No computation or regression script is a premise of the proof. The
exact-arithmetic and integer-programming checks are supplementary
audits. Using one consistent counting convention, they comprise 46,390
exact-arithmetic and LP checks and 91 exact ILP checks, for a total of
46,481 checks, with no reported discrepancy.

\subsection{\texorpdfstring{The public-preprint status line will be
updated only after all of the following are frozen: a
zero-\texttt{sorry} build, the axiom report, the Lean toolchain and
manifest, and the release commit. The intended final claim is therefore
already delimited, but it is not asserted in this
draft.}{The public-preprint status line will be updated only after all of the following are frozen: a zero-sorry build, the axiom report, the Lean toolchain and manifest, and the release commit. The intended final claim is therefore already delimited, but it is not asserted in this draft.}}\label{the-public-preprint-status-line-will-be-updated-only-after-all-of-the-following-are-frozen-a-zero-sorry-build-the-axiom-report-the-lean-toolchain-and-manifest-and-the-release-commit.-the-intended-final-claim-is-therefore-already-delimited-but-it-is-not-asserted-in-this-draft.}

\section{12. Potential Uses and Future
Directions}\label{potential-uses-and-future-directions}

\subsection{12.1 Universal linear integrality
gap}\label{universal-linear-integrality-gap}

\subsubsection{Problem 12.1}\label{problem-12.1}

Is there an absolute constant \(C\) such that

\[
\nu_3^*(G)-\nu_3(G)
\le C|V(G)|
\]

for every split graph \(G\)?

The present paper supplies supporting structures near the extremal
corridor but does not resolve the bulk regime without quadratic slack.

\subsection{12.2 Effective global
constants}\label{effective-global-constants}

\subsubsection{Problem 12.2}\label{problem-12.2}

Find an explicit global constant \(C\) in Theorem 1.1.

The near-extremal corridor is already effective by Proposition 10.1. A
fully quantitative global theorem therefore requires only explicit rates
in fractional-to-integral packing, or a structured replacement for
Haxell--Rödl in the bulk, together with effective dense-decomposition
thresholds in the sparse-independent regime.

\subsection{12.3 Algorithmic packing}\label{algorithmic-packing}

The proof contains polynomially implementable pieces:

\begin{itemize}
\tightlist
\item
  complete-graph factorizations;
\item
  weighted selection of reserved vertices;
\item
  bipartite list edge coloring;
\item
  matching extraction;
\item
  dense decomposition algorithms implicit in iterative absorption.
\end{itemize}

It is natural to seek a polynomial-time algorithm that outputs a clique
partition of size

\[
\frac{n^2}{6}+O(n)
\]

with an explicit constant. The factorization, list-coloring, and
matching steps of Sections 5--7 are individually polynomial; assembling
them into a single approximation algorithm with a proven guarantee is
left to future work.

\subsubsection{Corollary 12.2}\label{corollary-12.2-effective-corridor-algorithm}

In the near-extremal corridor covered by Proposition 10.1 (explicit
\(p_0,s_0\)), the proof is constructive: the one- and
double-factorizations of \(K_p\), the reserved-vertex selection, the
list edge coloring (Appendix D, whose kernel/stable-pairing proof is
itself an algorithm), and the matching extraction all run in polynomial
time. Hence there is a deterministic polynomial-time algorithm that, for
split graphs in the effective corridor, outputs a clique partition of
size \(n^2/6+2n\) with the explicit constant of Proposition 10.1.
Outside the corridor the guarantee currently relies on the
non-constructive bulk input (Theorem 2.1); a fully effective global
algorithm is left to future work.

\subsection{12.4 Stability and extremal
classification}\label{stability-and-extremal-classification}

\subsubsection{Problem 12.3}\label{problem-12.3}

Characterize split graphs satisfying

\[
\operatorname{cp}(G)
\ge
\frac{n^2}{6}-o(n^2).
\]

The proof suggests that near-extremal graphs should have
\(|I|=(2+o(1))|K|\) and low-dispersion absence profiles close to a
common center.

Two pieces of evidence support a Simonovits-type stability statement.
First, an exhaustive enumeration at \(n=9\) (all split graphs with
\(|K|=3\)) shows that graphs within an additive constant of the maximum
are within a bounded edit distance of the complete-split family, while
genuinely dispersed profiles fall a linear amount below the maximum.
Second, the fractional margin of Theorem 4.2 already yields a level-set
bound: away from near-extremality the quadratic slack \(\mu(\alpha)p^2\)
is strictly positive, so any near-maximizer must have \(\alpha\) near
the extremal ratio and small profile dispersion. A full stability
theorem would still require quantifying the gain of a single
symmetrization step and controlling the directions along which
\(\Phi^*\) is flat; we leave it as an open problem, most naturally
phrased in the cut (graphon) metric to absorb the choice of nearest
complete-split graph.

\subsection{12.5 Higher-clique linear-error
packing}\label{higher-clique-linear-error-packing}

The complete-split reduction determines the fractional extremal problem
and the first-order integral asymptotics for fixed \(K_r\)-and-edge
partitions. The natural higher-clique analogue of the present paper is
the following linear-error question.

For fixed \(r\ge3\), does every split graph satisfy

\[
\pi_r(G)
\le
\frac{r-1}{4r}n^2+O_r(n),
\]

where \(\pi_r(G)\) is the minimum number of parts in an edge partition
into copies of \(K_r\) and single edges?

The common-profile and cloning ideas suggest one possible approach, but
the relevant local orbit programs and integral rounding mechanisms
become substantially more complicated. Design and decomposition methods
such as those developed by Keevash {[}10{]} and the regular-slice
framework of Allen, Böttcher, Cooley, and Mycroft {[}11{]} may be useful
tools; they are not ingredients of the present proof.

\subsection{12.6 Chordal and clique-tree
applications}\label{chordal-and-clique-tree-applications}

Shifted-center completion was designed to separate owned clique edges
from protected interface edges. A parametric version may be useful in a
future chordal linear-error theorem, where separators are shared between
neighboring maximal cliques.

\subsection{12.7 Reusable proof
principle}\label{reusable-proof-principle}

The paper illustrates a broader strategy:

\begin{quote}
Use exact symmetric LPs to create slack away from extremality, then
reserve explicit combinatorial rounding only for the corridor where the
slack vanishes.
\end{quote}

This may apply to other structured packing-covering problems in which a
general regularity theorem is too weak at second order.

\begin{center}\rule{0.5\linewidth}{0.5pt}\end{center}

\section{13. Reproducibility}\label{reproducibility}

The working release package for this draft is organized as follows.

\begin{Shaded}
\begin{Highlighting}[]
\NormalTok{01\_MANUSCRIPT/}
\NormalTok{    PAPER\_III\_split\_lineal\_v0.9.4\_review\_en.md}
\NormalTok{02\_PROOF\_CONTROL/}
\NormalTok{    LEDGER.md}
\NormalTok{    LEDGER\_INCREMENTAL\_v0.9.2\_formalization\_draft.md}
\NormalTok{    AUDIT\_PAPER\_C.md}
\NormalTok{03\_FORMALIZATION/}
\NormalTok{    lean/                         (Lean 4 / Mathlib v4.28.0 sources)}
\NormalTok{    lakefile.toml}
\NormalTok{    lean{-}toolchain}
\NormalTok{    lake{-}manifest.json}
\NormalTok{    BUILD\_STATUS.md               (temporary until the release build is frozen)}
\NormalTok{04\_SUPPLEMENTARY\_REGRESSION/}
\NormalTok{    audit\_c\_fast.py}
\NormalTok{    audit\_c\_ilp.py}
\NormalTok{    audit\_c\_ilp\_results.txt}
\end{Highlighting}
\end{Shaded}

The manuscript is the mathematical source of truth. \texttt{LEDGER.md}
is the authoritative dependency specification for the Lean development;
the incremental ledger records only the transition from v0.9.1 to this
draft. The final release package will add the frozen commit, build log,
axiom report, and SHA-256 manifest.

The proof itself is contained in the manuscript. The common-profile
optimization is proved in Theorem 3.1; the fractional margin is proved
in Theorem 4.2 and Appendix A; the factorization, polarization, and
shifted-center inequalities are proved in Lemmas 5.1--7.1; the
divisibility correction is proved in Section 8.3 and Appendix B; and the
explicit near-extremal constants are proved algebraically in Sections
9.3 and 10.5.

\subsection{The included script is a supplementary regression test for
the explicit corridor constants. Its output is not used as a hypothesis,
reduction, or proof step anywhere in the manuscript. Historical
exploratory computations are not part of the logical proof
package.}\label{the-included-script-is-a-supplementary-regression-test-for-the-explicit-corridor-constants.-its-output-is-not-used-as-a-hypothesis-reduction-or-proof-step-anywhere-in-the-manuscript.-historical-exploratory-computations-are-not-part-of-the-logical-proof-package.}

\section{Appendix A. Algebra of the Fractional
Margin}\label{appendix-a.-algebra-of-the-fractional-margin}

Let \(x=d/p\) and \(\alpha=q/p\). Ignoring the exact linear terms
temporarily, the three normalized branches of \(F\) are

\[
f_1(x)=\frac16+\frac{\alpha x}{3},
\]

\[
f_2(x)=x^2-x+\frac12,
\]

and

\[
f_3(x)=\frac16+\frac{x^2}{3}.
\]

Subtracting the affine target contribution gives minima

\[
\min_x\left(f_1(x)-\frac{\alpha x}{2}-C_\alpha\right)
=
\frac{\alpha^2}{12},
\]

\[
\min_x\left(f_2(x)-\frac{\alpha x}{2}-C_\alpha\right)
=
\frac{(2-\alpha)^2}{48},
\]

and

\[
\min_x\left(f_3(x)-\frac{\alpha x}{2}-C_\alpha\right)
=
\begin{cases}
\alpha(8-5\alpha)/48,&\alpha\le4/3,\\
(2-\alpha)^2/12,&\alpha\ge4/3.
\end{cases}
\]

The third is dominated by the lower envelope of the first two. Restoring
the exact \(-p/2\) contribution in each pointwise estimate and averaging
produces the \(-p/4\) term in Theorem 4.2.

\begin{center}\rule{0.5\linewidth}{0.5pt}\end{center}

\section{Appendix B. Divisibility
Correction}\label{appendix-b.-divisibility-correction}

Let \(P=x_1\cdots x_p\) be a path and let \(O\subseteq V(P)\) have even
cardinality. Define

\[
J=\{x_jx_{j+1}:|O\cap\{x_1,\ldots,x_j\}|\text{ is odd}\}.
\]

Every internal vertex \(x_j\) changes prefix parity precisely when
\(x_j\in O\). Therefore

\[
\operatorname{Odd}(J)=O.
\]

Also \(|E(J)|\le p-1\) and \(\Delta(J)\le2\).

After parity correction, the minimum degree has fallen by at most two.
In the sparse-independent regime it is still greater than \(3p/4\) for
all sufficiently large \(p\), so Turán's theorem supplies a \(K_5\).
Deleting a \(C_4\) inside that \(K_5\) changes the edge count by
\(1\pmod3\), while deleting a \(C_5\) changes it by \(2\pmod3\). Every
affected degree changes by two, so parity is preserved. The total degree
loss from the path subgraph and the correcting cycle is at most four.
Thus, when \(q=o(p)\), the final graph has minimum degree at least
\(p-1-q-4\), which eventually exceeds \((0.9+\varepsilon_0)p\) for any
fixed \(0<\varepsilon_0<0.1\).

\begin{center}\rule{0.5\linewidth}{0.5pt}\end{center}

\section{Appendix C. Computational
Audits}\label{appendix-c.-computational-audits}

\subsection{C.1 Exact common-profile
LP}\label{c.1-exact-common-profile-lp}

\texttt{verify\_common\_profile\_lp.py} enumerates the vertices of the
symmetrized cover polyhedron using exact rational Gaussian elimination
and compares the optimum with (3.5).

\subsection{C.2 Exact fractional
margin}\label{c.2-exact-fractional-margin}

\texttt{verify\_fractional\_margin.py} checks Theorem 4.2 for

\[
3\le p\le80,
\qquad
1\le q\le2p,
\qquad
0\le d\le p
\]

using exact rational arithmetic.

\subsection{C.3 Small ILP audits}\label{c.3-small-ilp-audits}

\texttt{verify\_factor\_rounding.py} and
\texttt{verify\_shifted\_center.py} compute exact integral triangle
packings on small instances using a binary edge-capacity ILP. They
verify the stated upper bounds on \(\Phi(G)\).

\subsection{C.4 Polarization}\label{c.4-polarization}

\texttt{verify\_polarization.py} performs exhaustive small-profile and
randomized exact-integer checks of Lemma 6.1.

\subsection{C.5 Divisibility}\label{c.5-divisibility}

\texttt{verify\_divisibility\_correction.py} exhaustively verifies the
path parity construction through order eighteen, checks the fixed
minimum-degree threshold exactly, and follows the full parity and
modulo-three correction on randomized dense residuals.

\begin{center}\rule{0.5\linewidth}{0.5pt}\end{center}

\section{Appendix D. Self-Contained List Edge
Coloring}\label{appendix-d.-self-contained-list-edge-coloring}

This appendix proves the list edge coloring theorem used in Section 7.2
(stated as Theorem 2.2), so that the corridor argument depends on no
external coloring theorem. The result is the maximum-degree case of
Galvin's theorem {[}12{]}; the gain graph of Section 7.2 is simple and
bipartite, which is the case proved here.

\subsection{D.1 Kernels}\label{d.1-kernels}

A \textbf{kernel} of a digraph \(D\) is an independent set \(K\) of
vertices such that every vertex outside \(K\) has an out-neighbor in
\(K\). \(D\) is \textbf{kernel-perfect} if every induced subdigraph has
a kernel. A kernel of a nonempty digraph is nonempty.

\subsubsection{Lemma D.1}\label{lemma-d.1-kernel-coloring-lemma}

Let \(D\) be a kernel-perfect digraph in which every vertex \(v\) has a
list \(L(v)\) with \(|L(v)|\ge d^+_D(v)+1\). Then the underlying graph
of \(D\) has a proper coloring choosing each color from the
corresponding list.

\subsubsection{Proof}\label{proof-7}

Pick any color \(c\) appearing in some list, let \(S=\{v:c\in L(v)\}\),
and let \(K\) be a kernel of \(D[S]\). Color every vertex of \(K\) with
\(c\); this is proper on \(K\) because \(K\) is independent. Delete
\(K\) from \(D\) and delete \(c\) from the list of every vertex of
\(S\setminus K\). Every \(v\in S\setminus K\) lost one color but also at
least one out-neighbor, namely its out-neighbor in \(K\), so the
hypothesis \(|L(v)|\ge d^+(v)+1\) is preserved; vertices outside \(S\)
keep their lists while their out-degrees do not increase. Repeat. Every
uncolored vertex always has a nonempty list, and each round with
\(S\ne\varnothing\) colors the nonempty set \(K\), so every vertex is
eventually colored. \(\square\)

\subsection{D.2 Stable pairings}\label{d.2-stable-pairings}

Let \(B\) be a bipartite graph with parts \(U\) and \(R\) in which every
vertex \(z\) has a linear preference order \(>_z\) on its incident
edges. A matching \(M\subseteq E(B)\) is \textbf{stable} if every edge
\(f\notin M\) has an endpoint \(z\) that is covered by an edge
\(e\in M\) with \(e>_zf\).

\subsubsection{Lemma D.2}\label{lemma-d.2-galeshapley-13}

For every system of preferences, a stable matching exists.

\subsubsection{Proof}\label{proof-8}

Run deferred acceptance. While some \(u\in U\) is unmatched and has not
yet proposed along all of its edges, let \(u\) propose along its most
preferred untried edge, say \(f\) with endpoint \(r\in R\). If \(r\) is
unmatched, or prefers \(f\) to its currently held edge, then \(r\)
accepts \(f\), releasing its previous partner, if any; otherwise \(r\)
rejects \(f\). Every edge is proposed at most once, so the process
terminates, and it terminates in a matching \(M\).

Let \(f=ur\notin M\). Suppose first that \(u\) proposed along \(f\) at
some point. Then \(r\) either rejected \(f\) immediately or accepted it
and later released it; in both cases, at that moment \(r\) held or
received an edge it strictly prefers to \(f\). Since the edge held by
\(r\) only improves during the process, the final edge \(e\in M\) at
\(r\) satisfies \(e>_rf\). Suppose instead that \(u\) never proposed
along \(f\). An unmatched vertex of \(U\) keeps proposing while untried
edges remain, so at termination \(u\) is matched, say to \(e\in M\); and
since \(u\) proposes in decreasing order of preference and never reached
\(f\), we have \(e>_uf\). In both cases \(f\) has an endpoint whose
matched edge is preferred to \(f\). \(\square\)

\subsection{D.3 The coloring theorem}\label{d.3-the-coloring-theorem}

\subsubsection{Theorem D.3}\label{theorem-d.3-galvin-maximum-degree-case}

Let \(B\) be a simple bipartite graph with maximum degree \(\Delta\),
and let every edge \(e\) carry a list \(L(e)\) with \(|L(e)|\ge\Delta\).
Then \(B\) has a proper edge coloring choosing each color from the
corresponding list.

\subsubsection{Proof}\label{proof-9}

\textbf{Step 1: a proper \(\Delta\)-edge-coloring.} By König's edge
coloring theorem {[}14{]}, \(B\) has a proper edge coloring
\(\varphi:E(B)\to\{1,\ldots,\Delta\}\). For completeness: induct on the
number of edges. Remove an edge \(e=ur\) and color the rest. At most
\(\Delta-1\) colors appear at \(u\) and at most \(\Delta-1\) at \(r\),
so \(u\) misses a color \(\alpha\) and \(r\) misses a color \(\beta\).
If they miss a common color, give it to \(e\). Otherwise consider the
maximal path \(P\) starting at \(u\) whose edges are colored alternately
\(\beta,\alpha\). The path \(P\) cannot end at \(r\): since \(r\) misses
\(\beta\), such a path would end with an \(\alpha\)-edge and hence have
even length, while every path between the adjacent vertices \(u\) and
\(r\) of a bipartite graph has odd length. Swap the colors \(\alpha\)
and \(\beta\) along \(P\); by maximality of \(P\) the coloring stays
proper, the vertex \(r\) is untouched, and \(u\) now misses \(\beta\).
Give \(\beta\) to \(e\).

\textbf{Step 2: the orientation.} Define preferences from \(\varphi\):
every \(u\in U\) prefers its incident edges with \textbf{higher}
\(\varphi\)-color, every \(r\in R\) prefers \textbf{lower}
\(\varphi\)-color. Define a digraph \(D\) on the vertex set \(E(B)\):
for distinct edges \(e,f\) sharing an endpoint \(z\), put an arc
\(e\to f\) exactly when \(f>_ze\).

Out-degrees: let \(e=ur\) with \(\varphi(e)=c\). The out-neighbors of
\(e\) are the edges at \(u\) with color greater than \(c\) --- at most
\(\Delta-c\), since the colors at \(u\) are pairwise distinct --- and
the edges at \(r\) with color smaller than \(c\) --- at most \(c-1\).
Hence

\[
d^+_D(e)\le(\Delta-c)+(c-1)=\Delta-1.
\]

\textbf{Step 3: kernel-perfectness.} An induced subdigraph of \(D\) is
\(D[S]\) for a set \(S\) of edges of \(B\). A kernel of \(D[S]\) is
exactly a stable matching of the subgraph \((V(B),S)\) with the
inherited preferences: independence in \(D\) means that no two kernel
edges share an endpoint, because two edges at a common vertex are always
joined by an arc; and the domination condition asks that every
\(f\in S\setminus K\) have an arc into \(K\), that is, an endpoint \(z\)
and an edge \(e\in K\) at \(z\) with \(e>_zf\), which is stability.
Lemma D.2 provides a stable matching, so \(D\) is kernel-perfect.

\textbf{Step 4: conclusion.} Every edge has
\(|L(e)|\ge\Delta\ge d^+_D(e)+1\), so Lemma D.1 yields a proper coloring
of the underlying graph of \(D\) from the lists --- that is, a proper
list edge coloring of \(B\). \(\square\)

\subsubsection{Remark D.4}\label{remark-d.4-application}

The gain graph of Section 7.2 is simple and bipartite, and hypothesis
(7.2) bounds every list by \(b-t_i\ge\max\{\rho,u\}\ge\Delta\) of the
gain graph. Theorem D.3 therefore applies verbatim. The multigraph
version of Galvin's theorem holds as well {[}12{]} and is not needed
here.

\begin{center}\rule{0.5\linewidth}{0.5pt}\end{center}

\subsection{Acknowledgements}\label{acknowledgements}

The author is deeply grateful to his wife María Paz and to his children
Lucas, Juan Cristóbal, Francisca, Raimundo, and Benjamín for their love,
patience, and support.

\begin{center}\rule{0.5\linewidth}{0.5pt}\end{center}

\subsection{Use of AI-assisted tools}\label{use-of-ai-assisted-tools}

AI-assisted tools were used during exploratory, computational,
adversarial, organizational, and editorial stages, including systems
from Anthropic, Google, and OpenAI. They supported the testing of
candidate arguments, exact regression checks, proof organization, audit
preparation, and drafting. The author reviewed the mathematical content,
selected the final arguments, and remains solely responsible for the
claims, citations, code, and presentation. No AI system is listed as an
author.

\begin{center}\rule{0.5\linewidth}{0.5pt}\end{center}

\subsection{References}\label{references}

{[}1{]} B. Barber, D. Kühn, A. Lo, and D. Osthus, ``Edge-decompositions
of graphs with high minimum degree,'' \emph{Advances in Mathematics}
\textbf{288} (2016), 337--385.

{[}2{]} O. V. Borodin, A. V. Kostochka, and D. R. Woodall, ``List edge
and list total colourings of multigraphs,'' \emph{Journal of
Combinatorial Theory, Series B} \textbf{71} (1997), no. 2, 184--204.

{[}3{]} G.-T. Chen, P. Erdős, and E. T. Ordman, ``Clique partitions of
split graphs,'' in \emph{Combinatorics, Graph Theory, Algorithms and
Applications} (Beijing, 1993), World Scientific, 1994, pp.~21--30.

{[}4{]} G. A. Dirac, ``Some theorems on abstract graphs,''
\emph{Proceedings of the London Mathematical Society} (3) \textbf{2}
(1952), 69--81.

{[}5{]} F. Dross, ``Fractional triangle decompositions in graphs with
large minimum degree,'' \emph{SIAM Journal on Discrete Mathematics}
\textbf{30} (2016), no. 1, 36--42.

{[}6{]} P. Erdős, E. T. Ordman, and Y. Zalcstein, ``Clique partitions of
chordal graphs,'' \emph{Combinatorics, Probability and Computing}
\textbf{2} (1993), no. 4, 409--415.

{[}7{]} P. E. Haxell and V. Rödl, ``Integer and fractional packings in
dense graphs,'' \emph{Combinatorica} \textbf{21} (2001), 13--38.

{[}8{]} J. P. Traverso Gianini, ``Complete-split extremizers for a
fractional triangle-cover functional on chordal graphs'' (Paper II in
the series), preprint v1.0, July 2026.

{[}9{]} R. Yuster, ``Integer and fractional packing of families of
graphs,'' \emph{Random Structures \& Algorithms} \textbf{26} (2005),
110--118.

{[}10{]} P. Keevash, ``The existence of designs,'' arXiv:1401.3665,
revised 2018.

{[}11{]} P. Allen, J. Böttcher, O. Cooley, and R. Mycroft, ``Tight
cycles and regular slices in dense hypergraphs,'' \emph{Journal of
Combinatorial Theory, Series A} \textbf{149} (2017), 30--100.

{[}12{]} F. Galvin, ``The list chromatic index of a bipartite
multigraph,'' \emph{Journal of Combinatorial Theory, Series B}
\textbf{63} (1995), no. 1, 153--158.

{[}13{]} D. Gale and L. S. Shapley, ``College admissions and the
stability of marriage,'' \emph{The American Mathematical Monthly}
\textbf{69} (1962), no. 1, 9--15.

{[}14{]} D. König, ``Über Graphen und ihre Anwendung auf
Determinantentheorie und Mengenlehre,'' \emph{Mathematische Annalen}
\textbf{77} (1916), 453--465.

{[}15{]} J. P. Traverso Gianini, ``Affine profile reduction for
fractional triangle packings in split graphs'' (Paper I in the series),
preprint v1.0, July 2026.

\end{document}
